\documentclass[11pt]{amsart}
\usepackage[T1]{fontenc}
\usepackage[margin=1.05in]{geometry}
\DeclareFontShape{OT1}{cmr}{bx}{n}{<-> cmbx10}{}
\DeclareMathSizes{8}{8}{8}{8}
\DeclareMathSizes{9}{9}{8}{8}
\DeclareMathSizes{10.95}{10.95}{8}{8}
\newtheorem{theorem}{Theorem}
\newcommand{\OPT}{\operatorname{OPT}}
\newcommand{\Alg}{\operatorname{ALG}}
\title{The sharp approximation ratio for parent-restricted CP-net aggregation}
\date{November 27, 2026}
\begin{document}
\begin{abstract}
For each \(t\ge3\), the worst-case approximation ratio of Algorithm~1
of Ali, Yang and Zilles on \(t\) complete binary conditional preference
tables is exactly \(2-2/t\). The family
lifts the representative-voter example of Endriss and Grandi into
overlapping input parent sets. Re-optimizing over each such set
reproduces its input table. For four inputs the costs are \(6\) and \(4\),
which answers the source's four-thirds question in the negative.
A standard best-input average gives the matching upper bound.
\end{abstract}
\maketitle

\section{Swap model and algorithm}
For binary attributes \(X_1,\ldots,X_m,V\), a complete conditional
preference table (CPT) for \(V\) on a parent set \(P\subseteq
\{X_1,\ldots,X_m\}\) declares one strict preference in every assignment
to \(P\). Interpret its output bit as the preferred value of \(V\).
Every context \(x\in\{0,1\}^m\) specifies exactly one swap at \(V\),
between the two outcomes that differ only at \(V\). The table orders the
swap by its output on the restriction \(x|_P\). Write \(g_i(x)\) for the
bit chosen by the \(i\)-th input, with all contexts expanded. The cost
of an aggregate \(h\) is
\[
  F(h)=\sum_{x\in\{0,1\}^m}\sum_{i=1}^t
        [h(x)\ne g_i(x)] .
\]
Here brackets denote the indicator of a true statement. Unrestricted optimum \(\OPT\) allows any parent set. It takes a
pointwise majority and pays \(\sum_x k_x\), where \(k_x\) is the
minority vote count and a tie has \(k_x=t/2\). Each choice on each
context is realizable with the full parent set.

In Algorithm~1 of Ali, Yang and Zilles \cite[pp.~10437--10438]{ayz},
the allowed parent sets are subsets of one input parent set. For each
input parent set \(P\), it re-optimizes the table by pooling votes over
all contexts with the same \(x|_P\). It takes the best such candidate.
Ali, Yang and Zilles ask in their conclusion whether this rule always
achieves a \(4/3\)-approximation. The quantification is over this local
algorithm for the target attribute. The lower bound below is also
realizable by complete acyclic CP-nets with \(V\) as a sink and
identical unconditional preferences for the other attributes.

\section{Exact ratio}
\begin{theorem}
For every \(t\ge 3\), the worst-case ratio of Algorithm~1 on
\(t\) complete binary input CPTs is \(2-2/t\), evaluated over
instances of positive optimum. At \(t=4\) the ratio is \(3/2\).
The supremum over all input counts is \(2\).
\end{theorem}

\begin{proof}
For the upper bound, use the best-input argument underlying the
representative-voter rule of Endriss and Grandi \cite{eg}
and Theorem~4 of \cite{ayz}. Every input table is feasible for
Algorithm~1, so \(\Alg\le \min_i F(g_i)\). At a context with
\(k=k_x\), the sum of the costs of the \(t\) candidates is
\(2k(t-k)\), the number of ordered pairs of opposing votes. When
\(k>0\), it is at most \(2(t-1)k\). The same inequality holds for
\(k=0\). Averaging over the \(t\) candidates and summing,
\[
  \Alg \le \frac{1}{t}\sum_{i=1}^t F(g_i)
       \le \left(2-\frac{2}{t}\right) \OPT .
\]
In particular, zero optimum forces zero algorithm cost.

For sharpness take \(m=t\) and the complete tables
\[
       P_i=\{X_j:j\ne i\},\qquad
       g_i(x)=\bigvee_{j\ne i} x_j,\qquad i=1,\ldots,t .
\]
Every declared parent is relevant, as can be seen by setting all
other visible bits to zero. The entire profile is explicit with
\(2^{t-1}\) rows per table. At \(0^t\) all votes are zero. At each
unit vector \(e_i\), only \(g_i\) votes zero. With at least two
ones all votes are one. Thus the full OR is the unique unrestricted
optimal function for \(t\ge 3\) and pays exactly \(\OPT=t\).

Fix an input parent set \(P_i\). It pairs each context \(x\)
with \(x_i=0\) with the context formed by changing \(x_i\) to one.
For every assignment of the visible bits with at least one one,
both completions have strict-majority vote one. The unique
cost-minimizing output on every such pair is one. The all-zero
assignment of the visible bits instead pairs \(0^t\) and \(e_i\).
The two constant choices on this pair cost
\[
\begin{array}{c|cc|c}
\text{output on the pair} & 0^t & e_i & \text{sum}\\ \hline
0 & 0 & t-1 & t-1\\
1 & t & 1   & t+1
\end{array}
\]
so the unique choice is zero. The restricted optimum therefore
reproduces \(g_i\) on all contexts. Every CPT on a subset of \(P_i\)
is constant on the same pairs, so it cannot lower this cost.

Relative to full OR, \(g_i\) changes only \(e_i\). Its cost is
therefore \(t-1\) from the other unit contexts plus \(t-1\) from
\(e_i\), giving \(\Alg=2(t-1)\). Equivalently, \(g_i\) disagrees
with every different input \(g_j\) on exactly \(e_i,e_j\).
Every re-optimized candidate has this same cost. The ratio
is \(2(t-1)/t\), and these values approach \(2\) as \(t\) grows.
\end{proof}

\section{Worked case and prior results}
With \(t=4\) there are five attributes including \(V\), and each
input table has eight rows. If \(X_1\) is hidden, the all-zero
visible row pairs \(0000\) and \(1000\). Across the two contexts
there are three votes for one and five for zero. The algorithm chooses
zero there, paying three disagreements, and pays one on each of the
other three unit contexts. Full OR pays one on each of the four unit
contexts. Thus \(6/4=3/2>4/3\). The result concerns the stated
parent-set restriction. The full OR in this family has \(2^t\) rows,
twice the size of a single input table. The source already observes
that unrestricted optimization is efficient for constant \(t\)
\cite[Proposition~3]{ayz}.

The votes on the \(t\) unit contexts form the complement of the
identity-matrix profile of Endriss and Grandi
\cite[Example~2]{eg}; the other contexts are unanimous. That example
already gives \(2-2/t\) for representative-voter selection. In the
CPT transfer in \cite[Theorems~16--17]{ayz}, Algorithm~1 can recover
the unrestricted optimum even when selecting an input is nearly a
factor of two worse. Here the overlapping parent sets force each
restricted optimization to keep its input rule. This is the role of
the unanimous contexts, which constrain the context pairs in a way
that plain ballot distances do not record. An independent guarantee
of \(4/3\) for best-input aggregation of CPTs with at most one parent
also applies there to Algorithm~1 \cite[Theorem~5.1]{aoyz}. The
construction above has \(t-1\) relevant parents in each input.
Ali's dissertation restates the general question for the same rule,
renumbered Algorithm~3 \cite[Sec.~3.7]{aliThesis}.

\clearpage
\section*{Verification and technical disclosure}
The Python replay checks every context and restricted fiber in the
displayed family for \(3\le t\le12\). A partial Lean~4 certificate
checks the four-input costs and fiber, the strict \(4/3\) violation,
and an integer inequality in the upper-bound calculation. The
general CP-net argument is given in the paper. Internally reviewed,
not peer reviewed. Under human direction, ChatGPT Astra~6.2 (OpenAI)
generated the discovery proof and research packet. OpenAI Codex
independently checked the mathematics. GPT-6.2 Sol (OpenAI) checked
the literature and prepared the manuscript, exact replay, partial
Lean certificate and illustrations. DannyExperiments maintains the
repository and is the preferred citation identity. No AI system is
proposed as an author.

\begin{thebibliography}{9}
\raggedright
\bibitem{eg} U. Endriss and U. Grandi,
\emph{Binary aggregation by selection of the most representative voters},
AAAI 28(1) (2014), 668--674. doi:10.1609/aaai.v28i1.8798.
\bibitem{ayz} A. M. H. Ali, B. Yang and S. Zilles,
\emph{Approximation algorithms for preference aggregation using CP-nets},
AAAI 38(9) (2024), 10433--10441. doi:10.1609/aaai.v38i9.28911.
\bibitem{aoyz} A. M. H. Ali, D. Ogundare, B. Yang and S. Zilles,
\emph{Approximation ratio for preference aggregation using tree CP-nets},
AAMAS (2025), 78--86.
\bibitem{aliThesis} A. M. H. Ali,
\emph{Exact and approximate preference aggregation using conditional preference
networks}, PhD thesis, University of Regina, October 2025.
Handle: 10294/17240.
\end{thebibliography}

\end{document}
